\chapter{基于近似黎曼求解器的 SRHD-GSPH 模型}

\section{SRHD-GSPH 方法构造与变量定义}

本章采用自然单位制 $c=1$，考虑狭义相对论理想流体在实验室系中的拉格朗日形式
\begin{equation}
\frac{dD}{dt}=-D\nabla\cdot\boldsymbol{v},\qquad
\frac{d\boldsymbol{q}}{dt}=-\frac{1}{D}\nabla p,\qquad
\frac{d\hat{e}}{dt}=-\frac{1}{D}\nabla\cdot(p\boldsymbol{v}),
\end{equation}
其中
\begin{equation}
D=\Gamma\rho,\quad
\Gamma=(1-v^2)^{-1/2},\quad
\boldsymbol{q}=\hat{h}\Gamma\boldsymbol{v},\quad
\hat{e}=\hat{h}\Gamma-\frac{p}{D},\quad
\hat{h}=1+\varepsilon+\frac{p}{\rho}.
\end{equation}
为避免符号混用，本文统一约定：$D$ 为实验室系密度，$\rho$ 为静止系密度，$\hat{e}$ 为守恒能量变量。

原始变量到守恒变量的初始化转换为
\begin{equation}
D=\Gamma\rho,\quad
\boldsymbol{q}=\hat{h}\Gamma\boldsymbol{v},\quad
\hat{e}=\hat{h}\Gamma-\frac{p}{D},\quad
\hat{h}=1+\varepsilon+\frac{p}{\rho},\quad
p=(\gamma-1)\rho\varepsilon.
\end{equation}

\section{GSPH 离散格式与一阶界面构造}

本文采用一阶 GSPH（无二阶重构、无斜率限制器），粒子对离散更新式为
\begin{equation}
\dot{\boldsymbol{q}}_a
=-\sum_{b\in N(a)}M_b p_{ab}^{*}
\left(\frac{1}{D_a^2}\nabla W_{ab}(h_a)+\frac{1}{D_b^2}\nabla W_{ab}(h_b)\right),
\end{equation}
\begin{equation}
\dot{\hat{e}}_a
=-\sum_{b\in N(a)}M_b p_{ab}^{*}\boldsymbol{v}_{ab}^{*}\!\cdot\!
\left(\frac{1}{D_a^2}\nabla W_{ab}(h_a)+\frac{1}{D_b^2}\nabla W_{ab}(h_b)\right).
\end{equation}

对任意粒子对 $(a,b)$，定义法向
\begin{equation}
\boldsymbol{n}_{ab}=\frac{\boldsymbol{x}_a-\boldsymbol{x}_b}{|\boldsymbol{x}_a-\boldsymbol{x}_b|},\quad
u_R=\boldsymbol{v}_a\cdot\boldsymbol{n}_{ab},\quad
u_L=\boldsymbol{v}_b\cdot\boldsymbol{n}_{ab}.
\end{equation}
一阶左右状态直接取粒子值
\begin{equation}
Q_{L,ab}=(\rho_b,u_L,p_b)^T,\qquad
Q_{R,ab}=(\rho_a,u_R,p_a)^T.
\end{equation}
界面速度写为
\begin{equation}
\boldsymbol{v}_{ab}^{*}
=u_{ab}^{*}\boldsymbol{n}_{ab}
\left[\frac{\boldsymbol{v}_a+\boldsymbol{v}_b}{2}
-\frac{u_L+u_R}{2}\boldsymbol{n}_{ab}\right].
\end{equation}

\section{近似黎曼求解器（相对论波速）}

相对论特征速度取
\begin{equation}
\lambda_{\pm}(u,c_s)=\frac{u\pm c_s}{1\pm uc_s},\qquad
c_s^2=\frac{\gamma p}{\rho\hat{h}},\quad 0\le c_s<1.
\end{equation}

\subsection{Rusanov}
\begin{equation}
\lambda_{\max}=
\max\!\left(|\lambda_-(U_L)|,|\lambda_+(U_L)|,|\lambda_-(U_R)|,|\lambda_+(U_R)|\right),
\end{equation}
\begin{equation}
p^*=\frac12\!\left[p_L+p_R-\lambda_{\max}(q_R-q_L)\right],\quad
u^*=\frac{p_Lu_L+p_Ru_R-\lambda_{\max}(\hat{e}_R-\hat{e}_L)}{2p^*}.
\end{equation}

\subsection{HLL}
\begin{equation}
S_L=\min\{\lambda_-(U_L),\lambda_-(U_R)\},\qquad
S_R=\max\{\lambda_+(U_L),\lambda_+(U_R)\},
\end{equation}
\begin{equation}
p^*=\frac{S_Rp_L-S_Lp_R+S_LS_R(q_R-q_L)}{S_R-S_L},
\end{equation}
\begin{equation}
u^*=
\frac{S_Rp_Lu_L-S_Lp_Ru_R+S_LS_R(\hat{e}_R-\hat{e}_L)}
{(S_R-S_L)p^*}.
\end{equation}

\subsection{HLLC（闭合表达）}
取同样的 $S_L,S_R$，定义
\begin{equation}
m_K=D_K(S_K-u_K),\quad K=L,R,
\end{equation}
由动量跳跃条件可得接触波速
\begin{equation}
S^*=\frac{p_R-p_L+m_Lu_L-m_Ru_R}{m_L-m_R},
\end{equation}
中间压力
\begin{equation}
p^*=p_L+m_L(S^*-u_L)=p_R+m_R(S^*-u_R).
\end{equation}
实际实现中若出现 $m_L\approx m_R$ 或 $p^*<0$，退化到 HLL 以保证鲁棒性。

\section{一阶时间推进与原始变量恢复}

\subsection{守恒量推进}
\begin{equation}
\boldsymbol{q}_a^{n+1}=\boldsymbol{q}_a^n+\Delta t\,\dot{\boldsymbol{q}}_a,\qquad
\hat{e}_a^{n+1}=\hat{e}_a^n+\Delta t\,\dot{\hat{e}}_a.
\end{equation}

\subsection{由 $(\boldsymbol{q},\hat{e})$ 恢复 $\chi,\boldsymbol{v},\Gamma,\hat{h}$}
引入
\begin{equation}
\chi=\frac{p}{D},\qquad A=\hat{e}+\chi=\hat{h}\Gamma,\qquad \boldsymbol{q}=A\boldsymbol{v}.
\end{equation}
记 $q^2=\boldsymbol{q}\cdot\boldsymbol{q}$，则
\begin{equation}
f(\chi)=q^2+(\hat{e}+\chi)\!\left(\frac{\chi}{\gamma-1}-\hat{e}\right)
+\sqrt{(\hat{e}+\chi)^2-q^2}=0.
\end{equation}
Newton 迭代
\begin{equation}
\chi^{(m+1)}=\chi^{(m)}-\frac{f(\chi^{(m)})}{f'(\chi^{(m)})},
\end{equation}
\begin{equation}
f'(\chi)=\frac{2\chi+(2-\gamma)\hat{e}}{\gamma-1}
+\frac{\hat{e}+\chi}{\sqrt{(\hat{e}+\chi)^2-q^2}},
\end{equation}
并使用阻尼确保 $\chi>0,\ (\hat{e}+\chi)^2>q^2$。

恢复公式
\begin{equation}
\boldsymbol{v}=\frac{\boldsymbol{q}}{\hat{e}+\chi},\quad
\Gamma=\frac{\hat{e}+\chi}{\sqrt{(\hat{e}+\chi)^2-q^2}},\quad
\hat{h}=\sqrt{(\hat{e}+\chi)^2-q^2}.
\end{equation}

\subsection{$D,x,p,\rho,\varepsilon$ 更新}
位置推进采用前后时刻速度平均
\begin{equation}
\boldsymbol{x}_a^{n+1}=\boldsymbol{x}_a^n+\frac{\Delta t}{2}
\left(\boldsymbol{v}_a^n+\boldsymbol{v}_a^{n+1}\right).
\end{equation}
实验室系密度不由显式连续方程推进，而由核求和重算：
\begin{equation}
\widetilde{D}_a^{n+1}=\sum_b M_bW(\boldsymbol{x}_{ab}^{n+1},C_{\rm smooth}h_a^n),\ 
h_a^{n+1}=k\left(\frac{M_a}{\widetilde{D}_a^{n+1}}\right)^{1/d},\ 
D_a^{n+1}=\sum_bM_bW(\boldsymbol{x}_{ab}^{n+1},h_a^{n+1}).
\end{equation}
最后
\begin{equation}
p_a^{n+1}=\chi_a^{n+1}D_a^{n+1},\qquad
\rho_a^{n+1}=\frac{D_a^{n+1}}{\Gamma_a^{n+1}},\qquad
\varepsilon_a^{n+1}=\hat{h}_a^{n+1}-1-\frac{p_a^{n+1}}{\rho_a^{n+1}}.
\end{equation}

\subsection{CFL 时间步}
\begin{equation}
\Delta t=C_{\rm CFL}\min_a\frac{h_a}{\max_{b\in N(a)}\left(|\lambda_+|,|\lambda_-|\right)},
\quad \lambda_{\pm}=\frac{u_n\pm c_s}{1\pm u_nc_s},\quad 0<C_{\rm CFL}<1.
\end{equation}

\section{本章小结}

本章给出与代码实现一致的一阶 SRHD-GSPH 算法：一阶界面构造、相对论波速下的近似黎曼求解、守恒量推进、以及基于 $\chi$ 的原始变量恢复。关键修正点包括：统一 $D/\rho$ 定义、统一能量符号 $\hat{e}$、补全波速与 CFL 表达、并给出可实现的 HLLC 闭合及退化策略。由此可形成“可复现、可编程、可验证”的第四章算法主体。

